\section{Compatible sections on the whole floor}\label{sec:gluing}

We finish Theorem~\ref{thm:floor-generation} by constructing sections on
all strata at one common degree.  The construction follows the
pre-admissible and admissible section induction of Fujino
\cite[Section~4]{Fujino2000}, using the restriction criterion and finite
images proved in the preceding sections.  The point to retain is that
compatibility holds on every intersection before the sections descend to
the existing line on the reduced floor.

We proceed from lower strata to higher ones.  At each stage the restriction
criterion lets us choose all extensions of an already compatible lower
collection.  Those extensions need not be invariant under comparisons.
Finite products of their transports will impose invariance while raising
every prescribed lower section to the same power.  To make this possible,
we first show that transport preserves the prescribed lower restrictions,
including when a surface comparison contracts a floor curve.

For a positive degree \(k\) and \(d\in\{0,1,2,3\}\), a \emph{system of
tuples through dimension \(d\)} is a vector subspace
\[
 \mathcal V_d(k)\subset
 \bigoplus_{\dim Z\leq d} H^0(Z,L_Z^k)
\]
whose tuples satisfy the residue restriction equality along every
incidence among these strata.  We call the system invariant if, for each
\(e\in\{0,\ldots,\min(d,2)\}\), each tuple is compatible with every
arrow of \(\mathcal G_e\): the arrow carries its component at the source
to its component at the target.  It \emph{generates} if for every such
stratum \(Z\) and every \(z\in Z\), some tuple has its \(Z\)-component
nonzero at \(z\).  These are linear conditions except for generation.
The vector space is finite-dimensional because there are finitely many
strata and each is compact.

If a system generates, the span of the componentwise \(v\)th powers
of its tuples generates in degree \(kv\).  Restriction and transport
commute with products, so compatibility and invariance persist.  We may
therefore enlarge a successful degree to any sufficiently divisible later
degree.  Empty collections of strata impose no condition.

\begin{lemma}[Transport preserves lower restrictions]\label{lem:transport-lower-restrictions}
Fix \(k>0\) and \(d\in\{1,2\}\).  Let a tuple through dimension \(d\)
be compatible along all incidences, and suppose that its part through
dimension \(d-1\) is invariant.  For any arrow of \(\mathcal G_d\),
transport of the tuple's source component has the assigned lower
restriction on every floor stratum of its target.
\end{lemma}
\begin{proof}
For curves, the comparison is an isomorphism preserving the marked points
and their residue generators.  Invariance makes the lower point values
equal, so the assertion follows.

For surfaces, take a target floor curve and view its valuation on the
source of the comparison.  If its center is a floor curve there,
adjunction of the equality of meromorphic adjoint subsheaves on the graph
gives a crepant residue comparison of the two curve pairs.  This is an
arrow of \(\mathcal G_1\), and the assigned curve tuple is invariant.

If the center is a point, it is a zero-dimensional lc center of the
source dlt surface, by crepancy.  It is one of the point strata and is an
actual smooth crossing of two coefficient-one curves: the resolution
used in Proposition~\ref{prop:strata-adjunction} is an isomorphism at
such a point.  In coordinates \((x,z)\) for the crossing, a divisorial
lc place above it arises by successive blowups of crossings of two
coefficient-one branches.  To verify this, factor a surface resolution
into point blowups.  In an SNC surface boundary the new crepant
coefficient at the blowup of a crossing is the sum of the two coefficients
minus one; at a point on just one branch it is that coefficient minus one.
Starting with coefficients at most one, a new coefficient one can arise
only at a crossing of two coefficient-one branches.  This remains true
at every subsequent step.

The corresponding exceptional curve is a \(\PP^1\) whose different has
exactly the two adjacent coefficient-one points.  Residue of
\((dx/x\wedge dz/z)^{\otimes qk}\) along it is
\((dt/t)^{\otimes qk}\), up to a sign removed in the even degree: at each
crossing blowup the logarithmic change-of-coordinate determinant is
\(\pm1\).  Further point blowups do not change this meromorphic residue
on its strict transform.  The normal target floor curve is bimeromorphic,
hence isomorphic, to this \(\PP^1\), and its actual log line is trivial
with that generator.  Its section is determined by the common residue at
the two point strata.  Pulling a local surface section through the crossing
restricts on the exceptional curve to its value at the point times that
generator.  Compatibility of the original tuple identifies this value
with its assigned point component, and invariance makes it the assigned
value at both target point strata.  Thus the entire target curve receives
exactly the prescribed lower section.

The two alternatives apply to every composite surface comparison, so the
assertion holds for the whole groupoid, not only for a generating link.
\end{proof}

\begin{proposition}\label{prop:generating-tuples}
There is a degree \(k>0\) and a compatible generating system of tuples
through dimension three which is invariant in dimensions zero, one, and
two.
\end{proposition}
\begin{proof}
For point strata take the same scalar in the canonical zero-form generator
\(1\) on every point.  The even iterated residue convention identifies
these generators along all paths.  This gives an invariant generating
system in dimension zero.

Suppose a system has been constructed through dimension \(d-1\), for
\(1\leq d\leq3\).  Replace it by powers and their span in a sufficiently
large common degree, still denoted \(k\).  For a \(d\)-stratum \(Z\),
the components of any lower tuple glue to a section on its entire floor
\(C_Z\) by Proposition~\ref{prop:strata-adjunction}.  They satisfy its
link comparison: for \(d=1\) the link is a point comparison, for \(d=2\)
a curve comparison, and for \(d=3\) one of the generating surface links.
Lower invariance supplies each equality.  Proposition~\ref{prop:restriction-link}
therefore extends every such boundary section to \(Z\), once the degree
is sufficiently large.  There are only finitely many strata, so the degree
can be chosen uniformly.

Let \(\mathcal P_d(k)\) be the vector space of all tuples through
dimension \(d\) whose lower part belongs to the chosen lower system and
whose \(d\)-components restrict to those glued floor sections.  It maps
surjectively to the lower system, since the finitely many extensions can
be chosen independently.  This is a linear space of lifts; no choice of
a nonlinear extension operator is involved.

The pre-system \(\mathcal P_d(k)\) generates.  For \(z\in Z\) with
\(y=f_Z(z)\in f_Z(C_Z)\), choose a floor point above \(y\) and a lower
tuple nonzero there.  Any extension to \(Z\) is the pullback of a section
of \(N_Z^k\), by \eqref{eq:stratum-semiample-map}.  Its value at \(y\)
is nonzero, and hence it is nonzero at \(z\).  If \(y\notin f_Z(C_Z)\),
the sheaf \(\mathscr I_{f_Z(C_Z)}\otimes N_Z^k\) is generated for large
\(k\) by Serre's theorem.  A section nonzero at \(y\) pulls back to a
section vanishing on \(C_Z\); with zero lower tuple and all other new
components zero, it belongs to \(\mathcal P_d(k)\).  This also covers an
empty floor.  Lower points remain generated because the projection to
the lower system is surjective.

For \(d=3\) this pre-system is the required final system.  For \(d=1,2\)
we impose invariance by norm products.  Lemma~\ref{lem:transport-lower-restrictions}
applies to every pre-tuple: all transported top components have the same
assigned lower restrictions.  A product of \(h\) such factors will
therefore have the original lower tuple raised componentwise to \(h\).

For each orbit of \(d\)-strata choose a representative \(Z\), and let
\(G_Z\) be the finite isotropy image on \(H^0(Z,L_Z^k)\) from
Proposition~\ref{prop:finite-isotropy}.  Given a pre-tuple with component
\(s_Z\), form its norm
\[
 P_Z(s_Z)=\prod_{g\in G_Z}g(s_Z)
 \in H^0(Z,L_Z^{k|G_Z|}).
\]
Choose a common integer \(h\) divisible by every \(|G_Z|\), and use
\(P_Z(s_Z)^{h/|G_Z|}\).  Transport this section to every member of its
orbit.  This is well defined: a change of the chosen transport by an
isotropy arrow permutes the factors in degree \(k\), and transport
commutes with multiplication.  No finiteness statement about the
possibly larger representation in degree \(kh\) is needed.  Each factor
has the assigned lower restriction proved above, so the resulting lower
tuple is the original lower tuple raised componentwise to \(h\).
Thus all these sections are compatible across different orbits and all
deeper intersections, and they are invariant in dimension \(d\).

They still generate.  Fix a point \(z\) of an orbit member.  For each of
the finitely many transported factors, choose a point over \(z\) on a
resolution of its comparison graph.  Equality of the pulled-back
invertible adjoint subsheaves identifies the line fibers.  Nonvanishing
of that factor at \(z\) is therefore the nonvanishing of the representative
pre-section at a definite point of \(Z\).  Each such condition is a
nonzero linear evaluation functional on the generating vector space
\(\mathcal P_d(k)\).  A finite union of the proper kernels of these
functionals cannot cover a complex vector space.  One pre-tuple satisfies
all of them, and its norm is nonzero at \(z\).  At a lower point choose
a pre-tuple whose assigned lower value is nonzero; its \(h\)th power
remains nonzero.  Finally take the linear span of all the constructed
norm tuples.  Compatibility and invariance are linear conditions, so this
span is an invariant generating system in degree \(kh\).

This completes the induction through dimensions one and two; the
pre-system at dimension three then has all the required properties.
\end{proof}

\begin{proof}[Proof of Theorem~\ref{thm:floor-generation}]
Apply Proposition~\ref{prop:generating-tuples}.  Its three-dimensional
components agree on every subordinate stratum, so
Proposition~\ref{prop:strata-adjunction} descends each tuple uniquely to
a section of \(L^k|_S\).  For any \(x\in S\), choose a component through
\(x\).  The generating system has a tuple whose value there is nonzero.
It is the pullback of the descended value in the actual line fiber at
\(x\), so that descended value is nonzero.  Nakayama's lemma makes the
evaluation map onto \(L^k|_S\) surjective at \(x\).  The finite-dimensional
span of these global sections therefore generates at every point, proving
semiampleness on the whole reduced floor.
\end{proof}

For the supported model in Proposition~\ref{prop:supported-model}, choose
a common multiple of the actual adjoint index, the coefficients and actual
equivalence of \(P\), and the generated boundary degree just obtained.
It gives a line \(L=\OO_V(qA)\), a section with divisor \(qP\), and a
generated restriction on \(S\).  Its boundary sections define a holomorphic
map \(S\to\PP^b\) with the actual pullback identity for \(L|_S\).
The gluing argument has not extended these sections to \(V\); that is
the distinct lifting task to which we now turn.
